%قبل از شروع مسئله ترجیحاً چیزی ننویسید. اما اگر خواستید بنویسید، در قالب لیستی (فقرات و شمارش) بنویسید.
%برای اضافه کردن نمره، بعدی از شروع‌مسئله، [x] را اضافه کنید و به جای x  نمره مورد نظر را قرار دهید. برای زیرمسئله‌ها هم به همین صورت است.
%\شروع‌مسئله[10]

\شروع‌مسئله[10]
یک گراف \متن‌لاتین{n} راسی با یال‌های بدون جهت داریم. الگوریتمی با پیچیدگی زمانی \متن‌لاتین{O(n + m)} ارائه کنید و در آن یال های این گراف را طوری جهت‌دار کنید که کمترین تعداد مولفه‌ی قویا همبند تشکیل شود.

%اگر مسئله چند قسمت داشت به این صورت عمل کنید:
%\شروع{زیرمسائل}
%\زیرمسئله[0] % مسئلۀ ۱آ
%\زیرمسئله % مسئلۀ ۱ب
%\زیرمسئله % مسئلۀ ۱پ
%\زیرمسئله % مسئلۀ ۱ت
%\پایان{زیرمسائل}

\solution{
برای هر مولفه‌ی همبندی به طور جداگانه مساله را حل می‌کنیم چرا که واضح است رئوس دو مولفه‌ی جدا را نمی‌توان در یک مولفه قویا همبندی قرار داد.
ادعا می‌کنیم اگر گراف همبند باشد کمینه تعداد مولفه‌های قویا همبند که به آن می‌توان رسید برابر است با \متن‌لاتین{k + 1} که \متن‌لاتین{k} برابر است با تعداد یال‌های برشی گراف (یال برشی: یالی که با حذف آن تعداد مولفه‌های همبندی گراف افزایش می‌یابد). 
ابتدا ثابت می‌کنیم جواب حداقل برابر این تعداد است: \\
اگر یک یال برشی مثل یال \متن‌لاتین{(u, v)} داشته باشیم نمی‌توان کاری کرد که دو سر آن در یک مولفه‌ی قویا همبند قرار بگیرند چرا که با حذف این یال برشی، دوسر آن در دو مولفه‌‌ی همبندی جداگانه قرار می‌گیرند پس امکان ندارد طوری جهت دهی کنیم که هم \متن‌لاتین{u} به \متن‌لاتین{v} مسیر داشته باشد و هم بالعکس و تنها مسیر بین این دو راس، یال بین آن‌هاست درحالیکه باید حداقل دو مسیر بین این دو راس باشد تا هر دو به هم مسیر جهت‌دار داشته باشند.\\
بنابراین یال‌های برشی در جواب بهینه درون یک مولفه‌ی قویا همبند نیستند (چون ثابت کردیم دو سر آن نمی‌توانند در یک مولفه قویا همبند قرار بگیرند) و بین دو مولفه‌ی قویا همبند مختلف هستند پس با حذف آن‌ها تغییری در تعداد مولفه‌های قویا همبند در جواب بهینه ایجاد نمی‌شود. بنابراین می‌توانیم یال‌های برشی گراف را حذف کنیم و جواب ثابت می‌ماند. اگر \متن‌لاتین{k} یال برشی داشته باشیم بعد از حذف آن ها گراف به \متن‌لاتین{k + 1} مولفه‌ی همبندی تقسیم می‌شود که هرطور یال‌های آن را جهت‌دهی کنیم تعداد مولفه‌های قویا همبند حداقل \متن‌لاتین{k + 1} خواهد بود. بنابراین گراف اصلی هم که جوابش برابر جواب این گراف است حداقل به \متن‌لاتین{k + 1}  مولفه‌ی قویا همبند تقسیم می‌شود و از این کمتر نمی‌شود.\\
حال اثبات می‌کنیم می‌توان طوری یال‌ها را جهت‌دهی کرد که تعداد مولفه‌های قویا همبند دقیقا برابر \متن‌لاتین{k + 1} باشد: \\
روی گراف از یک راس \متن‌لاتین{dfs} می‌زنیم و یال‌های درخت را رو به پایین و \متن‌لاتین{back edge} ها را رو به بالا جهت‌دهی می‌کنیم. این کار را با \متن‌لاتین{O(n + m)} می‌توان انجام داد. به استقرای ریاضی ثابت می‌کنیم اگر یال‌های گراف را به این شکل جهت‌دهی کنیم تعداد مولفه‌های قویا همبند تشکیل شده برابر است با \متن‌لاتین{k + 1}.\\
لم: یال بین دو راس \متن‌لاتین{(v, p)} برشی  است اگر این یال در درخت \متن‌لاتین{dfs} آمده باشد و اگر \متن‌لاتین{p} پدر راس \متن‌لاتین{v} در درخت \متن‌لاتین{dfs} باشد، در زیردرخت راس \متن‌لاتین{v} هیچ یال \متن‌لاتین{back edge} به هیچ کدام از اجداد \متن‌لاتین{v} وجود نداشته باشد. اثبات آن هم واضح است چرا که اگر یالی در درخت نیامده باشد، با حذف آن یال، گراف همچنان با یال‌های زیردرخت \متن‌لاتین{dfs} همبند می‌ماند. اگر هم در درخت آمده باشد و شرط گفته شده را نداشته باشد، با حذف این یال دو راس \متن‌لاتین{v} و \متن‌لاتین{p} در دو مولفه همبندی جداگانه قرار می‌گیرند. برای اثبات این موضوع هم فرض خلف می‌کنیم که \متن‌لاتین{v} بعد از حذف این یال همچنان به راسی خارج از زیردرختش مثل \متن‌لاتین{p} راه داشته باشد. در این صورت مسیر از \متن‌لاتین{v} به \متن‌لاتین{p} را درنظر بگیرید. اولین جایی که به راسی خارج از زیر درخت \متن‌لاتین{v} می‌رویم از یک راس داخل زیر درخت \متن‌لاتین{v} به یک راس خارج از این زیردرخت رفته‌ایم درحالیکه فرض کردیم یالی از زیردرخت \متن‌لاتین{v} به خارج آن وجود ندارد که تناقض دارد. پس لم اثبات شد.\\
فرض استقرا: اگر گراف همبند ما \متن‌لاتین{k} یال برشی داشته باشد و یال های زیردرخت \متن‌لاتین{dfs} آن را رو به پایین و سایر یال‌ها را رو به بالا جهت‌دهی کنیم تعداد مولفه‌های قویا همبند تشکیل شده برابر است با \متن‌لاتین{k + 1}.\\
پایه: \متن‌لاتین{k = 0}\\
در این حالت هیچ یال برشی‌ای در گراف نداریم. می‌خواهیم ثابت کنیم گراف قویا همبند خواهد بود. از راس ریشه که از آن \متن‌لاتین{dfs} زده‌ایم که با یال‌های درخت (که رو به پایین جهت‌دهی شده‌اند) می‌توان به همه‌ی رئوس این مولفه رسید. حال ثابت می‌کنیم از هر راسی هم می‌توان به راس ریشه رسید. راس دلخواه \متن‌لاتین{u} را در نظر بگیرید. از آنجاییکه گراف طبق فرض یال برشی ندارد و با توجه به لم، حتما یک \متن‌لاتین{back edge} از زیردرخت \متن‌لاتین{u} به یکی از اجداد \متن‌لاتین{u} وجود دارد. از \متن‌لاتین{u} با یال‌های درخت به آن راس رفته (که همانطور که گفتیم رو به پایین جهت‌دهی شده‌اند) و سپس با \متن‌لاتین{back edge} گفته شده به جد \متن‌لاتین{u} می‌رویم (که جهت آن همانطور که گفتیم روبه بالا است). حال همین کار را برای جد \متن‌لاتین{u} که به آن رسیده‌ایم انجام می‌دهیم و به جد راس جدید می‌رویم. در هر مرحله با تکرار این کار به یکی از اجداد راسی که در آن هستیم می‌رویم و فاصله‌مان هربار با ریشه کمتر می‌شود تا اینکه بالاخره به ریشه می‌رسیم. پس با جهت‌دهی گفته شده از هر راسی به راس ریشه می‌توان رسید. پس گراف ما قویا همبند است چراکه راس ریشه به تمام رئوس راه دارد و تمام رئوس هم به ریشه راه دارند.\\
حکم: اگر گراف همبند ما \متن‌لاتین{k + 1} یال برشی داشته‌باشد و با روش گفته‌شده یال‌های آن را جهت‌دهی کنیم، تعداد مولفه‌های قویا همبندی برابر است با \متن‌لاتین{k + 2}.\\
چون 
$k + 1 \ge 1$
پس حتما یال برشی داریم. بین یال‌های برشی یالی که بیشترین عمق را دارد درنظر بگیرید (عمق: فاصله از ریشه). فرض کنید این یال بین راس \متن‌لاتین{p} و \متن‌لاتین{v} باشد که راس \متن‌لاتین{p} پدر راس \متن‌لاتین{v} است. زیردرخت راس \متن‌لاتین{v} را درنظر بگیرید. در این زیر درخت یال برشی نداریم چراکه اگر داشتیم عمق آن بیشتر از عمق یال \متن‌لاتین{(p, v)} بود و با این فرض که \متن‌لاتین{(p, v)} بین یال‌های برشی بیشترین عمق را دارد تناقض داشت. پس در زیر درخت راس \متن‌لاتین{v} هیچ یال برشی نداریم و یال‌های داخل درخت آن روبه پایین و سایر یال‌ها رو به بالا جهت‌دهی شده اند. پس طبق فرض استقرا که برای حالت پایه گفتیم (گراف بدون یال برشی) کل این زیردرخت یک مولفه‌ی قویا همبند خواهد بود. حال این زیر درخت را از گراف اصلی حذف می‌کنیم. گراف باقیمانده گرافی همبند با \متن‌لاتین{k} یال برشی است که طبق روش گفته شده جهت‌دهی شده است. پس بنابر فرض استقرا این گراف \متن‌لاتین{k + 1} مولفه‌ی قویا همبندی دارد که با یک مولفه‌ی قویا همبندی که حذف کردیم در کل \متن‌لاتین{k + 2} مولفه‌ی قویا همبندی داریم. 


}




%%%%%%%%%%%%%%%%%%%%%%%%%%%%%%%%%%%%%%%%%%%%%%%%%%%%%
% برخی دستورات مفید برای تمرین‌ها. %
%%%%%%%%%%%%%%%%%%%%%%%%%%%%%%%%%%%%%%%%%%%%%%%%%%%%%

% فرمول‌های نمونه:
%$f(x) = \Theta(x)$
%$T(x, y) \leq \log(x) + 2^y + \binom{2n}{n}$
% {\tt code\_function}


% لیست موارد بدون شماره‌گذاری:
%\شروع{فقرات}
%    \فقره مورد اول 
%        \شروع{فقرات}
%            \فقره مورد اول داخلی 
%                \شروع{فقرات}
%                    \فقره اولین مورد خیلی داخلی 
%                    \فقره دومین مورد خیلی داخلی 
%                \پایان{فقرات}
%            \فقره مورد دوم داخلی 
%        \پایان{فقرات}
%    \فقره مورد دوم 
%\پایان{فقرات}

% لیست موارد شماره‌گذاری شده:
%\شروع{شمارش}
%    \فقره یک مورد 
%    \فقره یک مورد دیگر 
%    \فقره باز هم یک مورد دیگر
%\پایان{شمارش}

% برای نوشتن یک فرمول چند خطی که یک شماره بخوره. برای حالت‌های دیگر به سایت زیر مراجعه کنید:
% https://www.overleaf.com/learn/latex/Aligning_equations_with_amsmath
%\begin{align} 
%    \begin{split}
%        (x+y)^3 &= (x+y)^2(x+y) \\
%                &= (x^2+2xy+y^2)(x+y) \\
%                &= (x^3+2x^2y+xy^2) + (x^2y+2xy^2+y^3) \\
%                &= x^3+3x^2y+3xy^2+y^3
%    \end{split}                                 
%\end{align}

% نوشتن ماتریس و آرایه دوبعدی
%\begin{align}
%    A = 
%    \begin{bmatrix}
%        A_{11} & A_{21} \\
%        A_{21} & A_{22}
%    \end{bmatrix}
%\end{align}

% اضافه کردن تصویر با شرح. برای اضافه کردن تصویر بدون شرح، فقط دستور درج‌تصویر را استفاده کنید.
%\شروع{شکل}[h!]
%\تنظیم‌ازوسط
%  \درج‌تصویر[پهنا=0.5\پهنای‌متن]{img.jpg}
%\شرح{شرح تصویر}
%\برچسب{تصویر:۱}
%\پایان{شکل}
%\پایان{زیرمسائل}

%کد زیر را ببینید.
%%\lstinputlisting{source_filename.py}
%\begin{latin}
%\begin{lstlisting}
%# Import the modules
%import sys
%import random
%
%ans = True
%
%while ans:
%question = raw_input("Ask the magic 8 ball a question: (press enter to quit) ")
%
%answers = random.randint(1,8)
%
%if question == "":
%sys.exit()
%
%elif answers == 1:
%print "It is certain"
%
%elif answers == 2:
%print "Outlook good"
%
%elif answers == 3:
%print "You may rely on it"
%
%elif answers == 4:
%print "Ask again later"
%
%elif answers == 5:
%print "Concentrate and ask again"
%
%elif answers == 6:
%print "Reply hazy, try again"
%
%elif answers == 7:
%print "My reply is no"
%
%elif answers == 8:
%print "My sources say no"
%\end{lstlisting}
%\end{latin}